\documentclass{article}
\usepackage{graphicx} % Required for inserting images
\usepackage{float}

\title{\textbf{Charger\\Progress log}}

\author{Nikhil,Madhur,Umesh}
\date{19 September 2026}

\begin{document}

\maketitle















\section{Day 1 [15 September 2026]}
We collected the equipment from the core labs. At first, we ran into minor issues, such as confusion regarding diode polarity conventions and polarized capacitors (filters). We used a multimeter, watched an instructional YouTube video on how to operate it, and verified that all components were functional.   
\begin{figure}[htbp]
    \centering
    \includegraphics[width=0.5\textwidth]{bread.jpeg}
   
    \label{fig:my_label}
\end{figure}
Once we understood the basics of the components, we started our work by building a bridge rectifier on a breadboard and testing it with the multimeter. The breadboard connection was successful. The primary challenge was soldering. After watching YouTube tutorials and observing other teams' techniques, we began experimenting. Our initial soldering attempt was poor, but we eventually got habituated the technique.




\begin{figure}[htbp]
    \centering
    \includegraphics[width=0.5\textwidth]{solder.jpeg}
   
    \label{fig:my_label}
\end{figure}
\newpage



We continuously verified connections as we added components—first checking the soldered bridge rectifier, then connecting it to the transformer, and testing along the way. Finally, the capacitor was soldered, and the voltage regulator was placed on the PCB without soldering.

\begin{figure}[htbp]
    \centering
    \includegraphics[width=0.5\textwidth]{Day1.jpeg}
   
    \label{fig:my_label}
\end{figure}







\section{Day 2 [17 September 2026]}
We returned to the core labs to complete the remaining soldering. We planned the layout for the voltage regulator, placing it away from other components because it generates significant heat. We initially doubted whether the regulator could handle converting 18 V DC down to 5 V, but confirmed via Google that it operates within an acceptable input range.

To ensure neat connections, we re-soldered a few components. Finally, we fixed the regulator to the PCB and verified the circuit with the multimeter. Attaching the USB port remains pending.

We also began 3D modeling for the enclosure. Madhur worked with FreeCAD while Nikhil tested OpenSCAD. We watched beginner tutorials for both programs, practicing basic 3D shapes to determine the best approach.






\section{Day 3 [18 September 2026]}
We attempted to solder the USB port directly to the PCB, but encountered difficulties because the pins were too small. We initially considered mounting the USB port vertically and taping it down. However, based on the TA's feedback regarding soldering quality, we decided to keep it horizontal. We chose to solder extension wires to the USB pins to make routing easier.


Regarding 3D modeling, we realized OpenSCAD was impractical for this task and switched fully to FreeCAD. After completing the introductory FreeCAD course, we learned basic operations (such as cutting and grouping) and created a initial 3D model.



\section{Day 4 [19 September 2026]}
We experienced technical issues with the soldering iron, as only the edge of the tip was heating properly. After trying to work around it, we postponed the remaining soldering to Monday. During this process, a transformer wire detached and some wire insulation burned, giving us extra repair work. As a result, the USB port could not be finalized.






\begin{figure}[htbp]
    \centering
    \includegraphics[width=0.5\textwidth]{Day4.jpeg}
   
    \label{fig:my_label}
\end{figure}

For the 3D model, we created an initial layout using estimated dimensions to figure out component placement. Instead of using screws, we decided to design small mounting pins (standoffs) that fit directly into the board holes. The initial draft of the 3D model was completed.



\begin{figure}[htbp]
    \centering
    \includegraphics[width=0.5\textwidth]{model1.jpeg}
    \label{fig:my_label}
\end{figure}




\section{Day 5 [21 September 2026]}
The USB port was successfully soldered, and the broken transformer wires were repaired. We tested the charger on a smartphone using the Ampere Flow app, but observed unstable wattage readings. We learned that shorting the D+ and D- data lines informs the phone to draw up to 5 W; otherwise, current draw is capped at 2.5 W. However, power output continued to fluctuate.

We measured the physical components using Vernier calipers. Realizing that our original 3D model did not fit the exact dimensions, we discarded the old design and began drafting a revised model.


\section{Day 6 [22 September 2026]}
After completing the TA session, we returned to our room for further testing. According to standard USB protocol, the phone should receive 2.5 W, but our circuit failed to deliver even 1 W (fluctuating between 0.2 W and 0.43 W). We began investigating the source of the power loss.   

\section{Day 7 [23 September 2026]}
We suspected issues with either the voltage regulator or the capacitor. We monitored the output using an oscilloscope and a USB power meter. The USB tester read approximately 4.68 V and 0.48 A. Checking the components revealed that the capacitor voltage was only 13.2 V (instead of the expected 17–18 V) and the regulator output matched this value, confirming both were malfunctioning.
\section{Day 8[24 September 2027]}
Before replacing the regulator, we tested the bridge rectifiers and discovered the diodes were damaged. We completely rebuilt the circuit using higher-capacity components rated up to 3 A, replacing all diodes, the capacitor, the regulator, and the USB port. We tested each stage of the rebuilt circuit successfully.



\begin{figure}[htbp]
    \centering
    \includegraphics[width=0.5\textwidth]{new.jpeg}
    
    \label{fig:my_label}
\end{figure}
\section{Day 9 [25 September 2026]}
We checked the power supplied to the phone (via the ampere flow app).\bigskip\\
\begin{figure}[htbp]
    \centering
    \includegraphics[width=0.5\textwidth]{aa.jpeg}
   
    \label{fig:my_label}
\end{figure}
\begin{enumerate}
    \item When the phone was connected to the circuit,the initial power draw was
2.90W(0.48A) and after that instant the power drops to 1.17W(0.33A).
    \item Running the test with other phones gave us wildly varying results.
\end{enumerate}
We were confused since the regulator supplies  5V approx(not exactly 5V but actually(4.98V), why do different phones have different voltage input and current flow ,that too the voltage fluctuation is observed even after it appeared to have stabilized. The charger was unable to increase the charge percentage but sometimes, it increased it (after 5 min for 1\%).\bigskip\\
    During testing we got a doubt that whether the
voltage drop after some time is due to the overheating of regulator,but confirmed that wasn't the reason (as soon as we remove and re-plug it repeats the same situation irrespective of the temperature of the regulator).\bigskip\\We suspect the fluctuation of the voltage is due to the inefficiency of the circuit(we have no concrete proof).







\end{document}


